\chapter{Trace Formula, Ordered Bar, and the Statistical Upside}
\label{v4-arith:tfe-gutzwiller}

\section{The Trace Formula Target}
\label{v4-arith:tfe:weil}

Arithmetic Koszul duality assembles an ordered-incidence
bar complex whose regularised character reproduces the Euler product of
\(\xi(s)\). The question it faces is: does this Euler-product data
already determine the zero distribution of \(\xi\), or does it leave
open a spectral problem? To answer precisely, fix the target.

\begin{definition}[Weil test-function class]
\label{v4-arith:tfe:defn:test-class}
\ClaimStatusDefinitional. Let \(\mathcal H\) be a standard Weil
test-function class: even Schwartz functions \(h : \mathbb{R} \to \mathbb{C}\)
whose Fourier transforms \(\widehat h(\tau) = \int_{-\infty}^\infty h(t) e^{i\tau t}\,dt\)
are compactly supported and sufficiently regular for the explicit
formula. The exact analytic class depends on the chosen normalization
of the Weil--Guinand formula.
\end{definition}

\begin{theorem}[Weil--Guinand explicit formula]
\label{v4-arith:tfe:thm:weil}
\ClaimStatusProvedElsewhere. Let \(\rho\) range over all nontrivial
zeros of the completed Riemann zeta function \(\xi(s)\), with no
critical-line parametrization assumed.
For \(h \in \mathcal H\), the completed zeta function satisfies the
distributional identity
\begin{equation}
\label{v4-arith:tfe:eq:weil}
\sum_\rho H(\rho)
=
W_\infty(h)+W_{\mathrm{pol}}(h)
-\sum_{p}\sum_{k\geq 1}
(\log p)p^{-k/2}
\bigl(\widehat h(k\log p)+\widehat h(-k\log p)\bigr),
\end{equation}
where \(H\) is the complex test function corresponding to \(h\), and
\(W_\infty(h)\) is the archimedean gamma-factor contribution
and \(W_{\mathrm{pol}}(h)\) is the contribution of the polar and
trivial zeros. Denote their sum by
\(W_{\mathrm{arch+pol}}(h) := W_\infty(h) + W_{\mathrm{pol}}(h)\).
\end{theorem}

\begin{remark}[three descriptions of the prime-power index set]
\label{v4-arith:tfe:rem:weil-sketch}
The geometric side of \eqref{v4-arith:tfe:eq:weil} is a sum over
prime powers \((p,k)\). In the Morishita dictionary this is a sum
over iterates \(K_p^k\) of the prime knot \(K_p\), with length
\(k\log p\). In a Gutzwiller language it is the periodic-orbit side.
In the ordered-bar language it is the logarithmic derivative of the
Euler-product character. These are three descriptions of the same
prime-power index set; none supplies the zero-side determinant-trace
operator.
\end{remark}

\section{Connes' Trace Formula}
\label{v4-arith:tfe:connes}

The Weil explicit formula is a distributional identity; the question
is whether the zero sum on its left side admits a determinant-trace
interpretation, i.e., whether some natural normal operator has
determinant divisor \(\mathrm{div}\,\xi_{\mathbb R}\) and
Weil--Connes trace supported at the ordinates \(\gamma_\rho\) after
zero-placement. Connes' adele-class-space construction provides a
partial answer.

\begin{theorem}[Connes--Meyer scaling trace]
\label{v4-arith:tfe:thm:connes-meyer}
\ClaimStatusProvedElsewhere. Connes' adele-class-space construction,
in Meyer's trace-formula formulation, realizes the Weil distribution
as a distributional scaling trace on the appropriate cokernel space
of the adele class space \(\mathbb{A}/\mathbb{Q}^*\):
\begin{equation}
\label{v4-arith:tfe:eq:connes-trace}
\operatorname{Tr}_{\mathrm{Connes}}(h)
=
\sum_\rho H(\rho)
\end{equation}
distributionally, where \(\rho\) ranges over the non-trivial zeros and
\(H\) is the Mellin transform paired with the test function \(h\). The
same prime-power and archimedean terms as \eqref{v4-arith:tfe:eq:weil}
occur. The scaling operator is not thereby a self-adjoint Hilbert-space
operator whose normal translate has determinant \(\xi(s)\); the
distributional trace is not by itself a spectral measure.
\end{theorem}

\section{Ordered-Bar Euler Skeleton}
\label{v4-arith:tfe:3d-acs}

The Connes--Meyer trace recovers the full Weil distribution
\eqref{v4-arith:tfe:eq:weil} as an operator trace. The ordered-incidence
bar complex approaches the same object from the Koszul-duality side:
its regularised character computes the Euler product. The question is
whether that character calculation closes the gap left by Connes.

\begin{proposition}[ordered-bar Euler skeleton]
\label{v4-arith:tfe:lem:ord-bar}
\ClaimStatusProvedHereConditional \emph{(on the ordered-incidence
model of \Cref{v4-arith:3d-acs:thm:chi-hom-bar} and on the chosen
regularised character)}. Let \(\chi_\infty(s)\) denote the
archimedean Euler factor arising from the regularisation at the
archimedean place of \(\overline{\mathrm{Spec}\,\mathbb{Z}}\). The
ordered-incidence bar complex
\(B^{\mathrm{ord}}_\bullet(\mathcal{V}^{\mathrm{prim}})\),
where \(\mathcal{V}^{\mathrm{prim}}\) is the Hochschild trace class
of the identity on the global arithmetic Iwahori--Hecke category of
\(GL_2\), has regularised character
\begin{equation}
\label{v4-arith:tfe:eq:ord-bar-product}
\operatorname{ch}_{\mathrm{reg}}
H_\bullet(B^{\mathrm{ord}}_\bullet(\mathcal V^{\mathrm{prim}}))
=
\chi_\infty(s)\prod_p(1-p^{-s})^{-1},
\end{equation}
with the completed \(\xi(s)\) obtained only after the archimedean and
polar normalisations are inserted.
\end{proposition}

\begin{proof}
This is the Euler-character computation of the ordered-incidence model:
finite prime factors contribute the Bost--Connes Euler factors, while
the archimedean and polar pieces are part of the regularisation data.
The statement is a character calculation, not a Lefschetz trace formula.
\end{proof}

\begin{remark}[canonical quantisation boundary]
\label{v4-arith:tfe:rem:canq}
If the profinite ACS theory and boundary comparison are assumed, the
same regularised character can be read as a boundary partition function.
That reading is conditional on
\Cref{v4-arith:3d-acs:conj:qu-acs} and
\Cref{v4-arith:3d-acs:thm:bulk-boundary-Vprim}.
\end{remark}

\section{Trace-Formula Equivalence}
\label{v4-arith:tfe:equiv}

The Euler-product character of Proposition~\ref{v4-arith:tfe:lem:ord-bar}
yields $\xi(s)$ only after regularisation. Under the regularised
character equality of \Cref{v4-arith:3d-acs:cor:partition-bar}, which
sets $Z(s) = \xi(s)$ in the ACS framework, the argument principle
converts the logarithmic derivative of $\xi$ into the zero sum; the
Weil--Guinand explicit formula then identifies this with the full Weil
distribution. The three trace functionals are conditionally equal.

\begin{theorem}[conditional equivalence of trace functionals]
\label{v4-arith:tfe:thm:equiv}
\ClaimStatusConditional \emph{(on
\Cref{v4-arith:3d-acs:conj:qu-acs},
\Cref{v4-arith:3d-acs:cor:BDCG-ordered-comparison}, and the
regularised character equality
\Cref{v4-arith:3d-acs:cor:partition-bar})}. Let \(\Gamma_T\) be a
positively oriented rectangle enclosing the zeros of \(\xi\) with
\(|\operatorname{Im}s|\leq T\), avoiding zeros on the boundary. The ACS argument-principle functional
\begin{equation}
\label{v4-arith:tfe:eq:z-pairing}
\mathcal Z_h:=
\frac{1}{2\pi i}\int_{\Gamma_T}
H(s)\frac{Z'(s)}{Z(s)}\,ds
\end{equation}
converges, after the standard \(T\to\infty\) limiting procedure and
the polar and trivial-zero subtractions, to the Weil explicit-formula
functional \eqref{v4-arith:tfe:eq:weil} whenever \(Z(s)=\xi(s)\) is
assumed.
\end{theorem}

\begin{proof}
Under the stated assumptions \(Z=\xi\). Applying the argument
principle to \(\xi'/\xi\) inside \(\Gamma_T\) gives the zero sum
\(\sum_{|\operatorname{Im}\rho|\leq T}H(\rho)\), together with polar
and trivial-zero residues. Expanding \(\xi'/\xi\) in the
half-plane \(\operatorname{Re}(s) > 1\) of absolute convergence
and performing the standard explicit-formula contour shift yields
the prime-power side of \eqref{v4-arith:tfe:eq:weil}. The equality
is inherited from the classical Weil--Guinand explicit formula.
\end{proof}

\begin{corollary}[equality of the ACS and Connes--Meyer distributions on \(\mathcal{H}\)]
\label{v4-arith:tfe:cor:info-equiv}
\ClaimStatusConditional \emph{(on the hypotheses of
\Cref{v4-arith:tfe:thm:equiv})}. For every \(h \in \mathcal H\), the
ACS functional \(\mathcal{Z}_h\) and the Connes--Meyer scaling trace
\(\operatorname{Tr}_{\mathrm{Connes}}(h)\) are equal.
\end{corollary}

\section{The Bar Reorganises, Not Locates}
\label{v4-arith:tfe:no-free}

The equivalence of the three trace functionals established conditionally
in Section~\ref{v4-arith:tfe:equiv} raises a sharper question: does the
ordered-bar construction supply an independent mechanism for placing the
zero ordinates $\gamma_n$ on the critical line, beyond what $\xi(s)$
itself encodes? The bar computation reorganises the zero data of $\xi$;
it does not add new data.

\begin{theorem}[bar computation reorganises zero data]
\label{v4-arith:tfe:thm:no-free-zeros}
\ClaimStatusProvedHere. The ordered-incidence bar complex determines
the individual zero ordinates \(\gamma_n\) only through the scalar
function \(\xi\): no additional self-adjoint operator or positivity
mechanism is produced by the bar calculation.
\end{theorem}

\begin{proof}
The bar supplies the Euler-product character and, after conditional
regularisation, the scalar function \(\xi\). The zero ordinates are
then the zero set of that scalar function. The bar calculation
produces no self-adjoint operator, spectral theorem, or positivity
input beyond what \(\xi\) already carries.
\end{proof}

\begin{remark}[the Hilbert--P\'{o}lya obstruction]
\label{v4-arith:tfe:rem:hilbert-polya}
A self-adjoint operator with spectrum $\{\gamma_n\}$ would place the
zeros on the critical line by the spectral theorem. Constructing such
an operator is precisely the spectral-flow obstruction of
\Cref{v4-arith:3d-acs:conj:spectral-flow}.
\end{remark}

\section{Arithmetic Gutzwiller}
\label{v4-arith:tfe:gutzwiller}

The zero ordinates $\gamma_n$ appear on the zero side of
\eqref{v4-arith:tfe:eq:weil} and, conditionally, as the spectrum of
the Hilbert--P\'{o}lya operator. The prime-power sum on the geometric
side of \eqref{v4-arith:tfe:eq:weil} takes the same form as the
periodic-orbit sum in the Gutzwiller trace formula: lengths $k\log p$
and amplitudes $(\log p) p^{-k/2}$. The prime-power side of the Weil formula is a periodic-orbit sum in
Gutzwiller's sense, conditional on the existence of an arithmetic
Hamiltonian whose primitive periodic orbits are the primes.

\begin{theorem}[classical Gutzwiller trace formula]
\label{v4-arith:tfe:thm:gutzwiller}
\ClaimStatusProvedElsewhere. Let $H$ be a quantum Hamiltonian whose
classical limit is chaotic with isolated primitive periodic orbits
$\gamma$. Write $d_{\mathrm{osc}}(E)$ for the oscillatory part of the
density of states $\sum_n \delta(E - E_n)$, $T_\gamma$ for the period
of orbit $\gamma$, $A_\gamma$ for its stability amplitude
(Maslov--Gutzwiller amplitude, incorporating the stability matrix
determinant), and $\mu_\gamma$ for its Maslov index. Then
\begin{equation}
\label{v4-arith:tfe:eq:gutzwiller}
d_{\mathrm{osc}}(E)\sim
\sum_\gamma A_\gamma
\exp(iET_\gamma-i\mu_\gamma\pi/2),
\end{equation}
asymptotically as $\hbar \to 0$.
\end{theorem}

\begin{theorem}[arithmetic Gutzwiller hypothesis]
\label{v4-arith:tfe:thm:arith-gutzwiller}
\ClaimStatusConditional \emph{(on a semiclassical ACS Hamiltonian,
an arithmetic phase space, and a prime-orbit correspondence)}. Call the
amplitudes \(A_p = (\log p)p^{-1/2}\) and lengths \(T_p = \log p\) the
\emph{Weil amplitudes} for the prime \(p\). If the ACS boundary
Hamiltonian has a classical chaotic limit whose primitive periodic
orbits are indexed by primes \(p\) with primitive period \(\log p\) and
stability amplitude \(A_p\), then the periodic-orbit expansion has the
desired form
\begin{equation}
\label{v4-arith:tfe:eq:arith-gutz}
\sum_p\sum_{k\geq 1}
(\log p)p^{-k/2}e^{itk\log p}
\quad+\quad\text{complex conjugate and archimedean terms}.
\end{equation}
\end{theorem}

\begin{proof}
This is substitution of the assumed arithmetic orbit dictionary into
the classical Gutzwiller formula. The proof is formal because the
semiclassical ACS phase space, the chaotic flow, and the amplitude
calculation are all hypotheses.
\end{proof}

\section{Connes--Meyer and the Gutzwiller Gap}
\label{v4-arith:tfe:connes-obstruction}

The arithmetic Gutzwiller hypothesis of the previous section requires a
smooth symplectic Hamiltonian system with a chaotic classical limit.
The question is whether the Connes--Meyer construction already supplies
this structure, or whether the Gutzwiller construction is genuinely new input.

\begin{theorem}[Connes--Meyer as distributional trace, not Gutzwiller system]
\label{v4-arith:tfe:thm:no-gutzwiller-connes}
\ClaimStatusProvedHere. The standard adele-class-space realization
used in the Connes--Meyer trace formula is a noncommutative scaling
trace construction. It is distinct from a smooth symplectic Hamiltonian
system of the kind required by the classical Gutzwiller theorem.
\end{theorem}

\begin{proof}
The adele class space \(\mathbb{A}/\mathbb{Q}^*\) is modeled by a
noncommutative quotient; its $C^*$-algebraic description gives an
operator-algebraic trace. The Gutzwiller formula requires a smooth
symplectic phase space, a Hamiltonian flow, isolated primitive periodic
orbits, and monodromy matrix determinants. None of these are part of
the Connes--Meyer construction as stated.
\end{proof}

\begin{remark}[other geometrisations remain open]
\label{v4-arith:tfe:rem:partial-geom}
Other geometrisations of the adele class space may supply a symplectic
structure; the theorem records only that the standard Connes--Meyer
trace formula is not itself a Gutzwiller mechanism. An arithmetic
Gutzwiller system, if it exists, is new structure beyond Connes--Meyer.
\end{remark}

\section{Statistical Upside}
\label{v4-arith:tfe:bgs}

The explicit formula \eqref{v4-arith:tfe:eq:weil} is a distributional
identity; it does not determine the individual zero ordinates $\gamma_n$
beyond what $\xi(s)$ already encodes. A Hilbert--P\'{o}lya operator, if
it exists, would put the determinant divisor on the critical line by
self-adjointness. The further statistical prediction, that the local
spacing statistics of the $\gamma_n$ follow the Gaussian Unitary
Ensemble (GUE), is not implied by either the explicit formula or
self-adjointness alone. It requires a dynamical mechanism: the
Bohigas--Giannoni--Schmit (BGS) principle connecting chaotic classical
dynamics to random-matrix statistics. The relevant symmetry class is
GUE because $\xi(s)$ has no anti-unitary time-reversal symmetry;
the unitary symmetry class applies. Let $R_2(u)$ denote the two-point
level spacing correlation function.

\begin{theorem}[Bohigas--Giannoni--Schmit principle]
\label{v4-arith:tfe:thm:bgs}
\ClaimStatusConjectured. Quantum systems whose classical limits are
chaotic are expected, after symmetry class is fixed, to exhibit random
matrix local spectral statistics. For the unitary symmetry class (GUE),
the two-point correlation function is
\begin{equation}
\label{v4-arith:tfe:eq:gue}
R_2(u)=1-\left(\frac{\sin \pi u}{\pi u}\right)^2.
\end{equation}
\end{theorem}

\begin{theorem}[conditional statistical upside]
\label{v4-arith:tfe:thm:stat-upside}
\ClaimStatusConditional \emph{(on
\Cref{v4-arith:3d-acs:conj:spectral-flow},
\Cref{v4-arith:tfe:thm:arith-gutzwiller},
arithmetic quantum chaos, and the BGS universality theorem
\Cref{v4-arith:tfe:thm:bgs})}. If the ACS Hamiltonian
of \Cref{v4-arith:3d-acs:conj:spectral-flow} has a chaotic
semiclassical limit with prime periodic orbits, then the local
statistics of the zero ordinates are governed by the corresponding
random-matrix ensemble; in the expected zeta case by GUE statistics
\eqref{v4-arith:tfe:eq:gue}.
\end{theorem}

\begin{proof}
Spectral flow supplies the zero-divisor determinant and, after
zero-placement, the Hamiltonian trace supported at the ordinates. Arithmetic
Gutzwiller identifies the classical periodic orbits with prime powers.
The BGS principle then transfers chaotic spectral statistics to the
zero ordinates. Every substantive step is one of the stated
hypotheses.
\end{proof}

\begin{remark}[GUE statistics as new content beyond the explicit formula]
\label{v4-arith:tfe:rem:content-added}
The Weil explicit formula establishes a distributional equality between
the zero ordinates and the prime-power sum; it implies neither
Montgomery pair correlation nor GUE spacing statistics. A 3D
Hamiltonian realization of the arithmetic ACS theory adds that
statistical mechanism, conditional on the semiclassical and
spectral-flow hypotheses of \Cref{v4-arith:3d-acs:conj:spectral-flow}.
\end{remark}

\begin{remark}[open conditions as explicit hypotheses]
\label{v4-arith:tfe:rem:residual}
\Cref{v4-arith:tfe:thm:stat-upside} rests on seven open conditions:
(1) the profinite quantum ACS on \(\overline{\mathrm{Spec}\,\mathbb{Z}}\);
(2) the Bezrukavnikov--Finkelberg and Costello--Gwilliam comparison
    with the ordered-incidence model;
(3) the regularised character equality \(Z=\xi\);
(4) the spectral-flow operator realizing Hilbert--P\'{o}lya;
(5) the full Weil--Guinand contour shift from the Euler skeleton;
(6) the arithmetic semiclassical phase space and Hamiltonian flow;
(7) chaos and random-matrix universality for that flow.
These are the explicit hypotheses of \Cref{v4-arith:tfe:thm:stat-upside}.
The ordered bar supplies the finite Euler skeleton
\eqref{v4-arith:tfe:eq:ord-bar-product}; the equality \(Z=\xi\) also
requires the archimedean, polar, and regularised character inputs.
Conditions~(1), (4), (6),
and~(7) constitute the spectral-flow problem of
\Cref{v4-arith:3d-acs:conj:spectral-flow}.
\end{remark}
